\documentclass[10pt,a4paper]{amsart}
\usepackage[margin=19mm]{geometry}
\usepackage{amsmath,amssymb,mathtools}
\usepackage[scr=rsfs,cal=euler]{mathalfa}
\usepackage{xfrac}
\usepackage[unicode,hidelinks,bookmarks=false]{hyperref}
\newcommand{\ie}{i.e.\ }
\newcommand{\cf}{cf.\ }
\newcommand{\resp}{respectively\ }
\newcommand{\Subsection}{\subsection}
\providecommand{\nobreakdash}{\nobreak\hbox{-}}
\setlength{\emergencystretch}{2em}
\multlinegap0pt
\usepackage{amssymb}
\usepackage{dsfont}
\usepackage{xcolor}
\usepackage{tikz}
\usepackage[shortlabels]{enumitem}
\newtheorem{claim}{Claim}
\def\cF{{\mathcal F}}
\def\fA{\mathfrak{A}}
\def\fB{\mathfrak{B}}
\title{The complete picture for clique factors in randomly perturbed graphs}
\author{Sylwia Antoniuk, Nina Kam\v{c}ev, Christian Reiher, Tadej Petar Tukara}
\date{Source: arXiv:2603.22081v1 (CC BY 4.0)}
\begin{document}
\maketitle

\noindent\emph{Source and licence.} The complete picture for clique factors in randomly perturbed graphs. Version arXiv:2603.22081v1 (CC BY 4.0), distributed by arXiv under the Creative Commons Attribution 4.0 International licence, http://creativecommons.org/licenses/by/4.0/, as recorded in the arXiv licence metadata for this exact version and shown on the abstract page https://arxiv.org/abs/2603.22081v1. Full licence notice: \texttt{licenses/arxiv-2603.22081-CC-BY-4.0.txt}.\par\smallskip

\noindent\textit{Editorial scope.} Counted unit: the article's claim c:p-B-edges (source0.tex lines 1177--1181). The article's complete original proof of it is delivered with it (source0.tex lines 1183--1236, one \texttt{proof} environment of 54 lines). No mathematics is written, repaired or re-derived: the statement and the proof are the article's own lines, in the article's own notation, and no retained line is substituted or re-flowed.  Retained uncounted as the source-local context the proof needs: lines 1110--1112; lines 1116--1119; lines 1140--1143.  Nothing was omitted from any retained span, so the package carries no omission record.  No retained line contains a citation, so the wrapper carries no bibliography and no bibliography entry.  No retained line contains a figure environment, an image-inclusion command or drawing code, so no image is delivered and the package declares no asset dependency.  Editorial formatting only, in addition: an AMS \texttt{amsart} preamble that restates the article's own declarations and macros used by the retained lines, namely the environments \texttt{claim} on separate counters, together with the standard packages those lines need.  The retained environments print in this document as Claim 1 in order of appearance, whereas the article prints the same items with the numbering of its own full document.  The complete native source of the article as distributed in the arXiv e-print for this exact version is bundled unchanged as \texttt{sources/196996/source0.tex}.
        \begin{equation}\label{eq:A_j}
            |A_j|\geq \frac{C}{m}.
        \end{equation}

        \begin{claim}
            \label{c:singular-A_m}
            If $t=0$, then $|A_m|< sm$ and hence $j<m$.
        \end{claim}

        \begin{claim}
            \label{c:p-A-edges-in-m+1}
            $$e(A_j,A_{m+1})\leq |A_j|\left((\alpha-\beta)|A_{m+1}|+C\right).$$
        \end{claim}

        \begin{claim}
            \label{c:p-B-edges}
            If $h\in[j-1]$, then 
            $$e(A_j,B_h)\leq |A_j|((\alpha-\beta)|B_h|+C).$$
        \end{claim}

        \begin{proof}
            The proof is similar to the proof of the previous claim. Again we assume for the sake of contradiction that 
            \begin{equation}\label{eq:B_h}
                |B_h|\geq C \quad \text{ and } \quad e(A_j,B_h)> (\alpha-\beta)|A_j||B_h|.
            \end{equation}
            Like in the proof of Claim~\ref{c:p-A-edges-in-m+1}, we construct an auxiliary bipartite graph $H$ between $\fA_j$ and $\fB_h$ by joining $\tilde K_j$ and $\tilde Q_h$ if the edge density between them is greater than $\alpha-2\beta$, that is, if 
            \[e(V(\tilde K_j), V(\tilde Q_h))> (\alpha-2\beta)j|V(\tilde Q_h)|.\]
            Notice that by the same reasoning as before, we have
            \[ e(H) \geq \beta|\fA_j||\fB_h|.\]
            Recall that $|V(Q_h)| = (m+1-h)r$. Fix an edge $f=\tilde K_j \tilde Q_h$ of $H$. 
            Let 
            \[
            X_f=\{x\in V(\tilde Q_h)\mid V(\tilde K_j) \subseteq N(x)\}.
            \]
            Note that 
            \[|X_f|  \geq e(V(\tilde K_j), V(\tilde Q_h))-(j-1)|V(\tilde Q_h)|\geq ((\alpha-2\beta)j-(j-1))(m+1-h)r.\] 
            We claim that, for sufficiently small $\beta$, the right-hand side in the second inequality above is larger than $(m-j)((m-h)s+t).$ Indeed, we have
            $$(\alpha j -j+1)(m+1-h)(ms+t)-(m-j)(ms-hs+t)= s(m-j)+t(j+1-h)>0,$$
            since we assumed that $h \leq j-1$, and since by Claim~\ref{c:singular-A_m} $j<m$ in case $t=0$.
            We conclude that
            \begin{equation}
             \label{eq:size-Xf}
            \begin{split}
            |X_f| %& \geq e(V(\tilde K_j), V(\tilde Q_h))-(j-1)|V(\tilde Q_h)|\\
            %&\geq j((m-1)s+t)(m+1-h)-(j-1)(m+1-h)(ms+t) \\
           % &=(m+1-h)((m-j)s+t)\\
            & > (m-j)((m-h)s+t).
            \end{split}
            \end{equation}

            If some vertex $x\in X_f$ is contained in the $L$- or $N$-set of $\tilde Q_h$, we can replace $\tilde K_j$ and $\tilde Q_h$ with a copy of $K_{j+1}$ (by shifting the vertex $x$ to $\tilde K_j$), $(m-h)s+t$ copies of $K_{m+1}$, one copy of $K_{h-1}$, and $s-t-1$ additional copies of $K_h$. This increases $\phi_{j+1}^\cF$ without decreasing $\phi_{m+1}^\cF$, which contradicts the maximality of $\cF$. For this reason $X_f$ is contained in the union of the $M$-sets of $\tilde Q_h$. 
            
            Now, by \eqref{eq:size-Xf} there is some index $i(f)\in [(m-h)s+t]$ such that $\tilde M_{i(f)}$ intersects $X_f$ in at least $m-j+1$ vertices. This means that there is a partition 
            \[\tilde M_{i(f)}=C_f\cup D_f\]
            with $|C_f|=m-j+1$ and $|D_f|=j-h$, and such that $C_f\subset X_f$. Since for each $\tilde Q_h$ the number of possible pairs $(i(f), C_f)$ as above is at most $((m-h)s+t){m+1-h\choose j-h} < \beta^{-1}$, we can choose one, say $(i(\tilde Q_h), C(\tilde Q_h))$, for which there are at least $\beta \deg_H(\tilde Q_h)$ edges $f$ with $(i(f), C_f) = (i(\tilde Q_h),C(\tilde Q_h))$.
            Thus, like in the proof of the previous claim, there is a~subgraph $H'$ of $H$ satisfying
            \begin{itemize}
                \item $e(H')\geq \beta e(H) \geq  \beta ^2 |\fA_j| |\fB_h|$,
                \item for any $\tilde Q_h$, there is a "special pair" $(\tilde M_i, \tilde N_i)$ in $\tilde Q_h$ for which there is a unique $\tilde C=C(\tilde Q_h)\subseteq \tilde M_i$, such that $(i,\tilde C)=(i(f),C_f)$ for any edge $f=\tilde K_j \tilde Q_h$ in $H'$. 
            \end{itemize}

            
            By \eqref{eq:A_j} we have $|\fA_j|\geq \frac{C}{m^2}$, and by \eqref{eq:B_h}, $|\fB_h|\geq \frac{|B_h|}{|Q_h|}\geq\frac{C}{(m-h)s+t}$. The K\H{o}vari-S\'os-Tur\'an theorem implies now that $H'$ contains a subgraph isomorphic to $K_{s-t,(m-j)s+t}$. 
            We can take this subgraph, which consists of $s-t$ copies of $K_j$ and $(m-j)s+t$ copies of $Q_h$, and replace them with one copy of $Q_j$, $((m-j)s+t)((m-h)s+t-1)$ copies of $K_{m+1}$ and some additional copies of $K_h$ as follows:
            \begin{itemize}
                \item the $s-t$ copies of $K_j$ become the $L$-part of $Q_j$,
                \item we take the $(m-j)s+t$ special pairs $(\tilde M_i, \tilde N_i)$, each from a different copy of $\tilde Q_h$, and split $\tilde M_i$ into $\tilde M_i = \tilde C \cup \tilde D$, where $\tilde C = C(\tilde Q_h)$; the sets $\tilde C$ become the $M$-part of $Q_j$ and the sets $\tilde D \cup \tilde N$ become the $N$-part of $Q_j$, 
                
                \item the remaining pairs $(\tilde M_k, \tilde N_k)$ are just copies of $K_{m+1}$, 
                \item finally, the $L$-part of $\tilde Q_h$ gives the additional copies of $K_h$.
            \end{itemize} 
            
            This increases $q_j^\cF$, while $\phi_{m+1}^\cF$ and $\phi_{j}^\cF$ remain unchanged, which is again a contradiction to the maximality of $\cF$.
        \end{proof}

\end{document}
